\chapter*{Resumen}

\section*{\tituloPortadaVal}

El acúmulo de contenidos en las plataformas digitales contemporáneas ha sobrepasado la capacidad
de los mecanismos clásicos de búsqueda y filtrado para acompañar al usuario en su exploración.
La búsqueda directa exige conocer de antemano lo que se busca, el filtrado por categorías opera
sobre taxonomías rígidas, y los sistemas de recomendación restan agencia al usuario al ocultar
los criterios sobre los que operan. Este Trabajo de Fin de Grado propone un motor de exploración
de colecciones digitales que combina tres mecanismos complementarios: filtrado facetado por
metadatos, navegación transversal entre colecciones relacionadas mediante referencias semánticas
y operaciones de conjuntos sobre los resultados parciales.

El sistema se articula como una arquitectura cliente-servidor. El núcleo de navegación se ha
desarrollado en Java sobre una base de datos relacional, exponiendo sus capacidades a través
de una \textit{API RESTful}. La interfaz web, desarrollada en JavaScript, da al usuario acceso
interactivo a los filtros, a los enlaces entre colecciones y al conjunto unión, que conserva
los resultados guardados durante la sesión. La validación funcional se ha realizado sobre
volcados de dos catálogos públicos de naturaleza distinta, TMDB y DBLP, y la usabilidad de la
interfaz se ha evaluado mediante un protocolo formal de pruebas con usuarios.


\section*{Palabras clave}

\noindent Navegación facetada, exploración de colecciones, hipermedia, filtrado relacional,
operaciones de conjuntos, motor de navegación, ingeniería del software, evaluación de
usabilidad, TMDB, DBLP.
